\chapter{INTRODUCTION}

% ============================================================================
% Section 1.1: Background of the Study
% ============================================================================
\section{Background of the Study}

Science education at the elementary level serves as a foundational stage in developing scientific literacy, conceptual reasoning, and problem-solving skills. Within the MATATAG Curriculum of the Department of Education (DepEd), Grade 6 Science introduces learners to fundamental concepts related to materials, including the description of changes of state for solids, liquids, and gases---specifically melting, evaporation, freezing, and condensation---using diagrams and flowcharts \cite{Kilag2024}. These competencies aim to help learners understand observable phenomena through scientific representations and structured reasoning.

Despite the importance of these competencies, research has consistently shown that students experience difficulty in understanding the particulate nature of matter and the mechanisms underlying phase changes. Learners often memorize definitions of melting or evaporation without fully comprehending the particle-level interactions that explain these transformations. Misconceptions frequently arise regarding particle spacing, motion, and the role of heat transfer in altering states of matter (Adadan, 2013; Paik, Kim, Cho, \& Park, 2004). These challenges are compounded by the abstract nature of microscopic representations, which are not directly observable.

In many classroom settings, demonstrations of phase change are limited to water-based experiments, such as melting ice or boiling water. While these activities provide observable macroscopic changes, they often lack explicit integration with particle-level models, resulting in fragmented conceptual understanding. Furthermore, access to laboratory materials for demonstrating processes such as sublimation or deposition is limited at the elementary level, restricting exposure to diverse examples of state change.

\begin{figure}[htbp]
    \centering
    \begin{subfigure}[b]{0.32\textwidth}
        \centering
        \includegraphics[width=\textwidth]{PhaseChangeDiagram.jpg}
        \caption{States of matter and phase changes}
        \label{fig:phase-change-diagram}
    \end{subfigure}
    \hfill
    \begin{subfigure}[b]{0.32\textwidth}
        \centering
        \includegraphics[width=\textwidth]{KhanAcademyPhaseChanges.png}
        \caption{Phase changes (Khan Academy)}
        \label{fig:khan-academy-phases}
    \end{subfigure}
    \hfill
    \begin{subfigure}[b]{0.32\textwidth}
        \centering
        \includegraphics[width=\textwidth]{88GuruPhaseChanges.png}
        \caption{Particle arrangement in phases (88Guru)}
        \label{fig:88guru-phases}
    \end{subfigure}
    \caption{Examples of phase change diagrams used in science education}
    \label{fig:phase-changes-examples}
\end{figure}

Game-Based Learning (GBL) has emerged as a promising instructional approach that integrates interactive systems with curricular objectives. Educational games can provide dynamic visualizations, immediate feedback, and experiential learning opportunities that reduce cognitive overload and promote conceptual change (Plass, Homer, \& Kinzer, 2015). When designed with explicit learning outcomes and aligned assessments, digital games may enhance engagement and deepen understanding, particularly in abstract scientific domains.

This study introduces \textit{Wiz's Journey: Mastering the States of Matter}, a 2D educational platformer developed as a supplementary and enrichment learning tool aligned with the MATATAG Grade 6 Science competency on changes of state. The game integrates explicit terminology (e.g., ``Solid $\rightarrow$ Liquid: Melting'') while allowing learners to discover appropriate phase changes through environmental problem-solving. Central to its design is an interactive particle-lattice minigame---a core mechanic of the gameplay---in which learners manipulate bonds and particle spacing to simulate state transitions. By directly engaging with representations similar to those used in classroom diagrams, learners repeatedly encounter visual models of particle behavior in a contextualized, goal-driven environment.

Given the increasing emphasis on learner-centered instruction under the MATATAG framework, there is a need to explore innovative digital tools that complement formal instruction without replacing it \cite{Abragan2022}. This study examines whether a platformer-based educational game can support conceptual understanding, engagement, and reinforcement of curriculum-aligned competencies related to phase changes of matter.


% ============================================================================
% Section 1.2: Objectives
% ============================================================================
\section{Objectives}

% Subsection 1.2.1: General Objective
\subsection{General Objective}

This study aims to design, develop, and evaluate an educational platformer game that rapidly operationalizes a selected learning competency from the first quarter of the Grade 6 MATATAG Science curriculum. The game will be developed with established pedagogical and assessment strategies in mind to ensure alignment with curriculum standards and learning outcomes. Learner performance and engagement will be evaluated through pretest and posttest analyses to determine the effectiveness of the intervention.

% Subsection 1.2.2: Specific Objectives
\subsection{Specific Objectives}

Specifically, the study aims to:

\begin{enumerate}
    \item Design learning outcome mechanics (LOMs) that strategically target both lower-order and higher-order cognitive learning outcomes aligned with the MATATAG Grade 6 Science competency on changes of state.
    \item Develop a functional 2D educational platformer game that incorporates the designed LOMs and integrates interactive particle-level simulations to visually represent state transitions.
    \item Evaluate the game's efficacy in teaching the selected learning outcomes to Grade 6 learners through structured pre- and post-assessment instruments, engagement surveys, and qualitative feedback from students regarding usability and perceived learning value.
\end{enumerate}


% ============================================================================
% Section 1.3: Scope and Limitations of the Study
% ============================================================================
\section{Scope and Limitations of the Study}

This study focuses exclusively on the development and evaluation of \textit{Wiz's Journey: Mastering the States of Matter}, a 2D educational platformer designed to reinforce the Grade 6 Science competency from the MATATAG Curriculum Guide (CG) 2023 for Grades 3-10. Specifically, the game targets the Quarter 1: Materials learning competency that requires students to ``describe changes of state for solids, liquids, and gases as melting, evaporation, freezing, and condensation using diagrams and flowcharts.'' The game addresses physical changes only and does not include instruction on chemical reactions or other chemical changes.

The game mechanics are designed to reflect the chosen learning outcomes, focusing on interactive methods for handling game content that directly map to the MATATAG competency. During the development phase, which commenced in December 2025 and extended through March 2026, the development team reviewed existing educational materials and pedagogical approaches to identify optimal mechanics and instructional strategies that could be contextualized for Grade 6 students' current knowledge of phase changes.

The game was developed using Unity as the game engine, chosen for its cross-platform compatibility and the development team's proficiency with the tool. The game is deployed on multiple platforms for greater accessibility: primarily on Android mobile devices via Itch.io, as well as on web platforms via Unity Play. This multi-platform approach ensures that the game can reach students regardless of their available device.

The evaluation methodology consisted of two phases: a formative playtest involving a small sample of students in a controlled environment with qualitative observations and surveys for iterative improvements, followed by a summative playtest with a larger sample that incorporated pretest, gameplay, and posttest assessments to measure learning outcomes.

However, recruitment and participation differed from the initially planned scope. The study originally planned to recruit participants from nearby schools including Inchican Elementary School, Seven Pillars Catholic School, and Princeton Academy. In actuality, only Inchican Elementary School was able to provide informed consent and participate in the study. The school provided 5 consenting students: 3 participants for the formative playtest and 18 participants for the summative playtest (compared to the initially promised 30 students). The other targeted schools were unable to participate due to logistical constraints, including lack of response to communication attempts and the timing of the study during the end of the academic school year when students were occupied with moving-up ceremonies from Grade 6 to first-year high school.

The study does not measure long-term retention of learning outcomes or academic performance beyond the duration of the intervention. Additionally, the research evaluates the game as a supplementary enrichment tool and does not employ a control group design, limiting the ability to make direct comparisons with traditional classroom instruction. The study adopts a single-group quasi-experimental design with pre- and post-assessments, engagement surveys, and qualitative interviews focused on student perception and usability feedback.


% ============================================================================
% Section 1.4: Significance of the Study
% ============================================================================
\section{Significance of the Study}

This study contributes to elementary science education in the Philippines by examining the integration of interactive digital modeling within a platformer game environment. For learners, \textit{Wiz's Journey: Mastering the States of Matter} provides repeated exposure to particle-level representations within meaningful problem-solving contexts, strengthening conceptual connections between diagrams and observable phenomena related to phase changes. By aligning the game with the MATATAG Grade 6 Science competency, the study demonstrates how educational games can reinforce curriculum-aligned learning objectives in accessible, engaging ways.

For educators and school administrators, this study offers insights into how educational games may complement MATATAG-aligned instruction without replacing formal teaching. The findings may inform future curriculum integration strategies, digital resource development, and professional development initiatives for teachers adopting game-based learning approaches. Given the resource constraints faced by many Philippine schools, this game's availability on multiple platforms (web and mobile) addresses accessibility barriers and demonstrates feasibility of deploying educational technology in resource-limited contexts.

For researchers and the broader educational community, this study adds to the growing body of literature on game-based learning in science education, particularly within the Philippine educational context and addressing the K-12 curriculum transition. By examining a discovery-based yet explicitly labeled game design combined with systematic evaluation using pre-post assessments and engagement surveys, the study contributes to ongoing discussions on balancing constructivist learning principles with guided instruction in digital environments. The qualitative feedback from student participants and actual implementation challenges encountered (such as participant recruitment constraints) provide realistic insights into deploying educational games in authentic school settings.
